\documentclass[tikz,border=4pt]{standalone}
\usepackage{tikz}
\usepackage{polymer-paper-preamble}
% Reusable Panel F apparatus motif. The caller owns labels, force arrows,
% panel framing, and whole-panel composition.
\newcommand{\PanelFFloatingCantilever}[1]{%
  \begin{scope}[shift={(#1)}]
    % Stable anchors for semantic overlays and downstream inspection.
    \coordinate (panelFFixedBoundary) at (0,0.34);
    \coordinate (panelFFloatingCantilever) at (-0.26,0);
    \coordinate (panelFDrivenElectrode) at (1.34,-1.40);
    \coordinate (panelFSourceReturn) at (1.74,0.39);
    \coordinate (panelFTrappedCharge) at (-0.22,-0.56);

    % Fixed mechanical boundary and shallow root jaw.
    \draw[cGray!62!black, line width=0.72pt, line cap=round, line join=round]
      (-0.26,0.34) -- (0.26,0.34);
    \foreach \hx in {-0.22,-0.08,0.06,0.20} {
      \draw[cGray!48!black, line width=0.30pt]
        (\hx,0.34) -- ({\hx+0.07},0.42);
    }
    \draw[cGray!62!black, line width=0.72pt, line cap=round]
      (0,0.34) -- (0,0.14);
    \fill[cGray!12, rounded corners=0.22mm]
      (-0.16,0) rectangle (0.16,0.14);
    \draw[cGray!62!black, line width=0.72pt, rounded corners=0.22mm]
      (-0.16,0) rectangle (0.16,0.14);
    \draw[cGray!56!black, line width=0.34pt] (-0.12,0) -- (-0.12,-0.06);
    \draw[cGray!56!black, line width=0.34pt] (0.12,0) -- (0.12,-0.06);

    % Electrically floating charge-trapping cantilever.
    \shade[top color=cAmber!18, bottom color=cAmber!46, rounded corners=0.5mm]
      (-0.14,0) .. controls (-0.02,-0.74) and (-0.49,-1.54) ..
      (-1.37,-2.04) -- (-1.22,-2.20) .. controls (-0.28,-1.72) and
      (0.24,-0.78) .. (0.14,0) -- cycle;
    \draw[cAmber!82!black, line width=0.62pt, rounded corners=0.5mm]
      (-0.14,0) .. controls (-0.02,-0.74) and (-0.49,-1.54) ..
      (-1.37,-2.04) -- (-1.22,-2.20) .. controls (-0.28,-1.72) and
      (0.24,-0.78) .. (0.14,0) -- cycle;
    \draw[cAmber!60!black, line width=0.28pt, opacity=0.42]
      (-0.06,-0.16) .. controls (-0.02,-0.76) and (-0.56,-1.54) .. (-1.29,-2.06);
    \draw[cAmber!60!black, line width=0.28pt, opacity=0.35]
      (0.04,-0.12) .. controls (0.14,-0.82) and (-0.39,-1.72) .. (-1.16,-2.16);

    % Driven electrode and gap guides.
    \shade[left color=cGray!30, right color=cGray!16]
      (1.34,-2.38) rectangle (1.58,-0.02);
    \draw[cGray!82!black, line width=0.62pt]
      (1.34,-2.38) rectangle (1.58,-0.02);
    \draw[cGray!45!black, line width=0.25pt] (1.39,-2.34) -- (1.39,-0.06);
    \foreach \hy in {-2.22,-1.98,-1.74,-1.50,-1.26,-1.02,-0.78,-0.54,-0.30} {
      \draw[cGray!48!black, line width=0.20pt]
        (1.58,\hy) -- (1.34,{\hy-0.05});
    }
    \foreach \yy/\xstart in {-1.88/-0.96,-1.50/-0.67,-1.12/-0.44,-0.74/-0.22} {
      \draw[cGray!23, line width=0.30pt, dash pattern=on 1.3pt off 1.2pt]
        (\xstart,\yy) -- (1.34,\yy);
    }
    \draw[cGray!19, line width=0.22pt, dash pattern=on 1.2pt off 1.5pt]
      (1.08,-0.42) -- (1.34,-0.42);

    % Driven source lead and grounded source return. No lead touches sample.
    \fill[cBlue!5, rounded corners=0.55mm] (0.76,0.40) rectangle (1.24,0.78);
    \draw[cBlue!64!black, line width=0.38pt, rounded corners=0.55mm]
      (0.76,0.40) rectangle (1.24,0.78);
    \node[font=\sffamily\bfseries\fontsize{4.4}{5.2}\selectfont,
          text=cBlue!70!black] at (1.00,0.59) {HV};
    \fill[cBlue!70!black] (1.24,0.50) circle (0.025);
    \draw[cBlue!70!black, line width=0.36pt, rounded corners=0.3mm]
      (1.24,0.50) -- (1.46,0.50) -- (1.46,-0.02);
    \fill[cBlue!70!black] (1.24,0.68) circle (0.025);
    \draw[cBlue!58!black, line width=0.34pt, rounded corners=0.3mm]
      (1.24,0.68) -- (1.74,0.68) -- (1.74,0.39);
    \draw[cBlue!58!black, line width=0.34pt] (1.61,0.39) -- (1.87,0.39);
    \draw[cBlue!58!black, line width=0.34pt] (1.65,0.33) -- (1.83,0.33);
    \draw[cBlue!58!black, line width=0.34pt] (1.69,0.27) -- (1.79,0.27);

    % Flat trapped-charge markers.
    \foreach \cx/\cy/\rr in {-0.22/-0.56/0.075,-0.41/-0.98/0.082,-0.74/-1.48/0.088,-1.12/-1.90/0.082} {
      \fill[cRed!68!black] (\cx,\cy) circle (\rr);
      \draw[cRed!88!black, line width=0.25pt] (\cx,\cy) circle (\rr);
    }
  \end{scope}%
}


\begin{document}
\begin{tikzpicture}[
  font=\sffamily,
  panel frame/.style={draw=cLGray, line width=0.32pt, rounded corners=1.2mm},
  panel letter/.style={font=\sffamily\bfseries\fontsize{9.5}{11.4}\selectfont, text=black},
  panel title/.style={font=\sffamily\bfseries\fontsize{7.2}{8.6}\selectfont, text=black, anchor=west},
  body label/.style={font=\sffamily\fontsize{5.5}{6.6}\selectfont, text=cGray},
  small label/.style={font=\sffamily\fontsize{4.8}{5.8}\selectfont, text=cGray},
  axis line/.style={-{Stealth[length=2.8pt,width=2.2pt]}, draw=cGray, line width=0.38pt},
  leader/.style={-{Stealth[length=2.5pt,width=1.9pt]}, draw=cGray, line width=0.32pt},
  electron/.style={circle, draw=cBlue!80!black, fill=cBlue!55, line width=0.30pt, minimum size=1.15mm, inner sep=0pt}
]

% Panel A
\begin{scope}[shift={(0,11.0)}]
  \draw[panel frame] (0,0) rectangle (7.2,5.0);
  \node[panel letter, anchor=west] at (0.28,4.62) {A};
  \node[panel title] at (0.82,4.62) {Sulfur-rich poly(S-r-DIB)};

  % --- elemental sulfur S8 crown (sulfur feedstock) ---
  \coordinate (s8c) at (1.16,3.92);
  \foreach \a/\i in {22.5/1,67.5/2,112.5/3,157.5/4,202.5/5,247.5/6,292.5/7,337.5/8} {
    \coordinate (v\i) at ($(s8c)+(\a:0.39)$);
  }
  \draw[cAmber!82!black, line width=0.5pt, line join=round]
    (v1)--(v2)--(v3)--(v4)--(v5)--(v6)--(v7)--(v8)--cycle;
  \foreach \i in {1,...,8} {
    \node[circle, draw=cAmber!85!black, fill=cAmberSphere!72, line width=0.30pt,
          minimum size=1.48mm, inner sep=0pt,
          font=\sffamily\bfseries\fontsize{3.1}{3.8}\selectfont, text=cBrown] at (v\i) {S};
  }
  \node[small label] at (1.86,3.92) {S$_8$};

  % --- inverse-vulcanization transformation ---
  \draw[-{Stealth[length=3.0pt,width=2.4pt]}, cGray, line width=0.5pt]
    (2.10,3.66) .. controls (2.82,3.46) and (3.12,3.20) .. (3.40,2.92);
  \node[small label, text=cBrown!92!black] at (3.02,4.02) {inverse vulcanization};
  \node[small label] at (2.72,3.16) {$\Delta$};

  % --- repeat unit: [ -S_x- meta-DIB (bis-thiocumyl) ]_n ---
  \node[body label] at (1.95,1.44) {sulfur segment};
  % left polysulfide -S-S-S-
  \foreach \x in {1.52,1.94,2.36} {
    \node[circle, draw=cAmber!85!black, fill=cAmberSphere!72, line width=0.42pt,
          minimum size=2.6mm, inner sep=0pt,
          font=\sffamily\bfseries\fontsize{4.9}{5.9}\selectfont, text=cBrown] at (\x,2.30) {S};
  }
  \draw[cAmber!85!black, line width=0.60pt] (1.69,2.30)--(1.77,2.30) (2.11,2.30)--(2.19,2.30);
  \draw[cGray, line width=0.58pt] (2.53,2.30)--(2.86,2.30);
  % left DIB-derived benzylic carbon (bare methyl stub = line-angle CH3)
  \coordinate (bcl) at (3.02,2.30);
  \draw[cGray, line width=0.50pt] (bcl)--($(bcl)+(0,-0.34)$);
  \draw[cGray, line width=0.58pt] (bcl)--(3.49,2.45);

  % DIB benzene ring (meta 1,3), Kekule aromatic (inner double bonds offset inward)
  \coordinate (p1) at (4.05,2.45);
  \foreach \a in {0,60,120,180,240,300} { \coordinate (r\a) at ($(p1)+(\a:0.56)$); }
  \draw[cGray, line width=0.62pt, line join=round] (r0)--(r60)--(r120)--(r180)--(r240)--(r300)--cycle;
  \draw[cGray, line width=0.44pt] ($(r0)!0.2!(r60)!0.14!(p1)$)--($(r0)!0.8!(r60)!0.14!(p1)$);
  \draw[cGray, line width=0.44pt] ($(r120)!0.2!(r180)!0.14!(p1)$)--($(r120)!0.8!(r180)!0.14!(p1)$);
  \draw[cGray, line width=0.44pt] ($(r240)!0.2!(r300)!0.14!(p1)$)--($(r240)!0.8!(r300)!0.14!(p1)$);
  \node[body label] at (4.05,3.42) {DIB-derived segment};

  % right DIB-derived benzylic carbon (from r300) + continuing polysulfide (horizontal)
  \coordinate (bcr) at ($(r300)+(0.34,0)$);
  \draw[cGray, line width=0.58pt] (r300)--(bcr);
  \draw[cGray, line width=0.50pt] (bcr)--($(bcr)+(0,-0.34)$);
  \draw[cGray, line width=0.58pt] (bcr)--($(bcr)+(0.24,0)$);
  \coordinate (sc2) at ($(bcr)+(0.48,0)$);
  \foreach \x in {0,1} {
    \node[circle, draw=cAmber!85!black, fill=cAmberSphere!72, line width=0.40pt,
          minimum size=2.3mm, inner sep=0pt,
          font=\sffamily\bfseries\fontsize{4.4}{5.3}\selectfont, text=cBrown] at ($(sc2)+(0.38*\x,0)$) {S};
  }
  \draw[cAmber!85!black, line width=0.58pt] ($(sc2)+(0.165,0)$)--($(sc2)+(0.215,0)$);
  \draw[cGray, line width=0.56pt] ($(sc2)+(0.50,0)$)--(5.92,1.965);

  % repeat-unit brackets [ ... ]_n
  \draw[cGray, line width=0.52pt] (1.10,1.20)--(0.84,1.20)--(0.84,3.02)--(1.10,3.02);
  \draw[cGray, line width=0.54pt] (0.84,2.30)--(1.26,2.30);
  \draw[cGray, line width=0.52pt] (5.60,1.20)--(5.86,1.20)--(5.86,3.02)--(5.60,3.02);
  \node[font=\sffamily\fontsize{6.2}{7.4}\selectfont, text=cGray] at (6.02,1.32) {$n$};
  \node[body label, align=center] at (3.45,0.56) {poly(S-\emph{r}-DIB) linear repeat unit};
\end{scope}

% Panel B
\begin{scope}[shift={(7.55,11.0)}]
  \draw[panel frame] (0,0) rectangle (7.2,5.0);
  \node[panel letter, anchor=west] at (0.28,4.62) {B};
  \node[panel title] at (0.82,4.62) {Composition series};

  % increasing sulfur content encoded as polysulfide chain length (S60 -> S85)
  \draw[-{Stealth[length=3.4pt,width=2.6pt]}, cAmber!90!black, line width=0.8pt]
    (0.66,1.45)--(0.66,3.85);
  \node[small label, rotate=90, anchor=south, text=cBrown] at (0.44,2.65) {increasing sulfur content};
  \foreach \y/\lbl/\sc in {1.55/S$_{60}$/2, 2.62/S$_{75}$/4, 3.70/S$_{85}$/6}{
    \coordinate (dl) at (1.62,\y);
    \draw[cGray, line width=0.5pt, line join=round]
      ([shift=(0:0.20)]dl)--([shift=(60:0.20)]dl)--([shift=(120:0.20)]dl)--([shift=(180:0.20)]dl)--([shift=(240:0.20)]dl)--([shift=(300:0.20)]dl)--cycle;
    \draw[cAmber!80!black, line width=0.55pt] (1.82,\y)--({2.24+0.36*\sc},\y);
    \foreach \k in {1,...,\sc}{
      \node[circle, draw=cAmber!85!black, fill=cAmberSphere!78, line width=0.32pt,
            minimum size=2.1mm, inner sep=0pt,
            font=\sffamily\bfseries\fontsize{3.9}{4.7}\selectfont, text=cBrown]
            at ({2.06+0.36*\k},\y) {S};
    }
    \coordinate (dr) at ({2.44+0.36*\sc},\y);
    \draw[cGray, line width=0.5pt, line join=round]
      ([shift=(0:0.20)]dr)--([shift=(60:0.20)]dr)--([shift=(120:0.20)]dr)--([shift=(180:0.20)]dr)--([shift=(240:0.20)]dr)--([shift=(300:0.20)]dr)--cycle;
    \node[body label, anchor=west, text=cBrown] at ({2.44+0.36*\sc+0.34},\y) {\lbl};
  }
  \node[small label, align=center] at (3.9,0.80) {DIB unit linked by polysulfide of increasing sulfur rank};
\end{scope}

% Panel C
\begin{scope}[shift={(0,4.75)}]
  \draw[draw=cAmber!70!black, line width=0.72pt, rounded corners=1.4mm] (0,0) rectangle (14.75,5.85);
  \node[panel letter, anchor=west] at (0.28,5.47) {C};
  \node[panel title] at (0.82,5.47) {Localized shallow and deep traps};
  \node[font=\sffamily\bfseries\fontsize{5.0}{6.0}\selectfont, text=cAmber!90!black,
        anchor=east] at (14.42,5.47) {HERO};

  % Real-space view: both trap classes are localized sites in one disordered film,
  % not separate material phases.
  \node[body label, font=\sffamily\bfseries\fontsize{5.8}{7.0}\selectfont] at (3.42,4.88) {real space};
  \fill[black, opacity=0.12, rounded corners=1.2mm] (0.91,1.13) rectangle (6.30,4.21);
  \shade[top color=cAmber!8, bottom color=cAmber!28, rounded corners=1.2mm]
    (0.82,1.22) rectangle (6.21,4.30);
  \draw[cAmber!70!black, line width=0.58pt, rounded corners=1.2mm]
    (0.82,1.22) rectangle (6.21,4.30);
  \draw[white, opacity=0.62, line width=0.38pt] (0.98,4.23)--(6.05,4.23);
  \foreach \yy/\phase in {3.63/0,2.78/18,1.92/-12} {
    \draw[cBrown!72!black, line width=0.56pt]
      (1.12,\yy) .. controls (1.65,{\yy+0.34}) and (2.14,{\yy-0.30}) .. (2.74,\yy)
      .. controls (3.32,{\yy+0.31}) and (3.92,{\yy-0.28}) .. (4.48,\yy)
      .. controls (5.02,{\yy+0.28}) and (5.55,{\yy-0.24}) .. (5.91,\yy);
    \foreach \xx in {1.42,2.70,4.45,5.62} {
      \shade[ball color=cAmberSphere!80] (\xx,{\yy+0.055*cos(\phase+\xx*70)}) circle (0.045);
    }
  }
  \foreach \x/\y in {1.72/3.48,4.88/2.92} {
    \shade[ball color=cBlue!66] (\x,\y) circle (0.095);
    \draw[cBlue!85!black, line width=0.26pt] (\x,\y) circle (0.095);
  }
  \foreach \x/\y in {2.56/1.78,5.26/1.62} {
    \shade[ball color=cRed!68] (\x,\y) circle (0.11);
    \draw[cRed!88!black, line width=0.26pt] (\x,\y) circle (0.11);
  }
  \node[small label, text=cBlue!85!black, anchor=west] at (1.02,0.78) {shallow sites};
  \draw[leader, cBlue!78!black] (2.02,0.86)--(1.74,3.33);
  \node[small label, text=cRed!88!black, anchor=east] at (6.02,0.78) {deep sites};
  \draw[leader, cRed!82!black] (5.12,0.86)--(5.25,1.47);
  \node[small label, text=cBrown, anchor=north east] at (6.03,4.12) {poly(S-\emph{r}-DIB) film};

  % Energy-space view: continuous shallow/deep trap DOS below the mobility edge.
  \draw[cLGray, line width=0.32pt] (6.72,0.52)--(6.72,4.88);
  \node[body label, font=\sffamily\bfseries\fontsize{5.8}{7.0}\selectfont] at (10.78,4.88) {energy diagram};
  \draw[axis line] (7.48,0.72)--(7.48,4.52);
  \node[body label, rotate=90] at (7.10,2.58) {energy};

  \draw[cGray!72, dash pattern=on 3pt off 2pt, line width=0.36pt] (7.72,4.25)--(12.98,4.25);
  \node[small label, anchor=west] at (13.10,4.25) {vacuum};
  \draw[cGray!88!black, line width=0.52pt] (7.72,3.78)--(12.98,3.78);
  \node[body label, anchor=west] at (13.10,3.78) {$E_C$};
  \draw[cGray!70, dash pattern=on 3pt off 2pt, line width=0.38pt]
    (7.72,3.30)--(8.05,3.30) (9.58,3.30)--(12.98,3.30);
  \node[small label, anchor=south west] at (8.05,3.36) {mobility edge};
  \draw[cGray!88!black, line width=0.52pt] (7.72,0.86)--(12.98,0.86);
  \node[body label, anchor=west] at (13.10,0.86) {$E_V$};

  \fill[cBlue!5, rounded corners=1.2mm] (7.80,2.47) rectangle (13.30,3.22);
  \fill[cRed!4, rounded corners=1.2mm] (7.80,1.02) rectangle (13.30,2.32);
  \shade[left color=cBlue!48, right color=cBlue!8]
    (7.82,2.52) .. controls (8.68,2.62) and (9.10,2.82) .. (9.28,2.91)
                 .. controls (9.08,3.02) and (8.62,3.16) .. (7.82,3.20) -- cycle;
  \draw[cBlue!82!black, line width=0.66pt]
    (7.82,2.52) .. controls (8.68,2.62) and (9.10,2.82) .. (9.28,2.91)
                 .. controls (9.08,3.02) and (8.62,3.16) .. (7.82,3.20);
  \shade[left color=cRed!45, right color=cRed!8]
    (7.82,1.08) .. controls (8.64,1.18) and (9.42,1.52) .. (9.68,1.70)
                 .. controls (9.42,1.92) and (8.62,2.22) .. (7.82,2.28) -- cycle;
  \draw[cRed!82!black, line width=0.72pt]
    (7.82,1.08) .. controls (8.64,1.18) and (9.42,1.52) .. (9.68,1.70)
                 .. controls (9.42,1.92) and (8.62,2.22) .. (7.82,2.28);

  \foreach \y in {2.72,3.04} {
    \draw[cBlue!82!black, line width=0.62pt] (9.18,\y)--(10.15,\y);
    \shade[ball color=cBlue!70] (9.62,\y) circle (0.085);
  }
  \foreach \y in {1.46,1.88} {
    \draw[cRed!82!black, line width=0.68pt] (9.42,\y)--(10.38,\y);
    \shade[ball color=cRed!72] (9.88,\y) circle (0.09);
  }
  \node[body label, font=\sffamily\bfseries\fontsize{6.0}{7.2}\selectfont,
        text=cBlue!82!black, anchor=west] at (10.62,2.90) {shallow};
  \node[body label, font=\sffamily\bfseries\fontsize{6.0}{7.2}\selectfont,
        text=cRed!82!black, anchor=west] at (10.62,1.68) {deep};

  \draw[-{Stealth[length=2.8pt,width=2.1pt]}, cBlue!78!black, line width=0.44pt]
    (9.60,3.04) .. controls (9.43,3.22) and (9.58,3.43) .. (9.73,3.63);
  \draw[-{Stealth[length=2.8pt,width=2.1pt]}, cRed!76!black, line width=0.44pt, dashed]
    (9.88,1.88) .. controls (10.12,2.30) and (10.10,2.78) .. (9.95,3.56);
  \node[small label, anchor=west, text=cGray!85!black] at (10.30,3.48) {thermal escape};
  \draw[{Stealth[length=2.7pt,width=2.0pt]}-{Stealth[length=2.7pt,width=2.0pt]},
        cRed!72!black, line width=0.42pt] (13.72,1.10)--(13.72,3.20);
  \node[body label, rotate=90, text=cRed!80!black] at (14.08,2.15) {$\Delta E_t$};
\end{scope}

% Panel D
\begin{scope}[shift={(0,0)}]
  \draw[panel frame] (0,0) rectangle (4.75,4.35);
  \node[panel letter, anchor=west] at (0.28,3.97) {D};
  \node[panel title] at (0.82,3.97) {Transient current};

  % Constant-voltage transient-current measurement across a MIM stack.
  \fill[cLGray!20, rounded corners=0.5mm] (0.45,2.48) rectangle (4.30,3.51);
  \draw[cGray, line width=0.38pt, rounded corners=0.5mm] (0.45,2.48) rectangle (4.30,3.51);
  \node[draw=cGray!80!black, fill=white, line width=0.36pt, rounded corners=0.45mm,
        minimum width=0.78cm, minimum height=0.56cm,
        font=\sffamily\bfseries\fontsize{4.7}{5.7}\selectfont, text=cGray!85!black] (smu) at (0.98,2.99) {SMU};
  \draw[cGray!82!black, fill=cGray!22, line width=0.38pt] (1.82,3.16) rectangle (3.66,3.30);
  \draw[cAmber!82!black, fill=cAmber!24, line width=0.42pt] (1.82,2.75) rectangle (3.66,3.16);
  \draw[cGray!82!black, fill=cGray!22, line width=0.38pt] (1.82,2.61) rectangle (3.66,2.75);
  \foreach \x in {1.96,2.24,...,3.52} {
    \draw[cGray!65, line width=0.26pt] (\x,3.16)--({\x+0.10},3.30);
    \draw[cGray!65, line width=0.26pt] (\x,2.61)--({\x+0.10},2.75);
  }
  \node[small label, text=cBrown] at (2.74,2.96) {polymer film};
  \draw[cGray, line width=0.38pt] (smu.east)--(1.82,2.99);
  \draw[cGray, line width=0.38pt] (3.66,2.68)--(4.03,2.68)--(4.03,2.57);
  \draw[cGray, line width=0.34pt] (3.90,2.57)--(4.16,2.57) (3.94,2.51)--(4.12,2.51);
  \node[small label, anchor=east] at (4.18,3.67) {MIM stack};

  % Symbolic log--log response: distributed traps give a power law; Debye is the contrast.
  \node[body label, anchor=west] at (0.72,2.27) {$I(t)\sim t^{-n}$};
  \draw[axis line] (0.66,0.47)--(0.66,2.10);
  \draw[axis line] (0.66,0.47)--(4.28,0.47);
  \node[body label, rotate=90] at (0.28,1.28) {$\log I$};
  \node[body label, anchor=north east] at (4.26,0.34) {$\log t$};
  \foreach \y in {0.84,1.26,1.68} {\draw[cGray!65, line width=0.32pt] (0.66,\y)--(0.75,\y);}
  \foreach \x in {1.46,2.38,3.30} {\draw[cGray!65, line width=0.32pt] (\x,0.47)--(\x,0.56);}

  \draw[cBlue!82!black, line width=0.78pt] (0.86,1.88)--(4.02,1.12);
  \draw[cRed!82!black, line width=0.82pt] (0.86,1.88)--(4.02,0.66);
  \foreach \x/\yb/\yr in {1.28/1.78/1.72,2.22/1.55/1.36,3.25/1.30/0.96} {
    \shade[ball color=cBlue!68] (\x,\yb) circle (0.052);
    \shade[ball color=cRed!68] (\x,\yr) circle (0.052);
  }
  \draw[cGray!88!black, dash pattern=on 3pt off 2pt, line width=0.52pt]
    (0.86,1.88)--(2.04,1.88)--(2.34,0.56);
  \node[small label, font=\sffamily\bfseries\fontsize{5.0}{6.0}\selectfont,
        text=cBlue!82!black, rotate=-14] at (3.26,1.53) {low $n$};
  \node[small label, font=\sffamily\bfseries\fontsize{5.0}{6.0}\selectfont,
        text=cRed!82!black, rotate=-22] at (3.34,0.78) {high $n$};
  \node[small label, text=cGray!88!black, anchor=south west] at (2.48,0.55) {Debye};
\end{scope}

% Panel E
\begin{scope}[shift={(5.00,0)}]
  \draw[panel frame] (0,0) rectangle (4.75,4.35);
  \node[panel letter, anchor=west] at (0.28,3.97) {E};
  \node[panel title] at (0.82,3.97) {ISPD trap distribution};

  % Non-contact induced electrostatic-voltmeter measurement (not a Kelvin probe).
  \fill[cLGray!18, rounded corners=0.5mm] (0.40,2.64) rectangle (4.35,3.53);
  \draw[cGray, line width=0.38pt, rounded corners=0.5mm] (0.40,2.64) rectangle (4.35,3.53);
  \draw[cGray!82!black, fill=cGray!20, line width=0.40pt] (1.18,2.77) rectangle (2.92,2.91);
  \draw[cAmber!82!black, fill=cAmber!24, line width=0.38pt] (1.38,2.91) rectangle (2.76,3.05);
  \foreach \x in {1.58,1.98,2.38} {
    \node[font=\sffamily\bfseries\fontsize{4.8}{5.8}\selectfont, text=cRed!76!black] at (\x,3.16) {$+$};
  }
  \draw[cGray!80!black, line width=0.42pt] (2.06,3.46)--(2.06,3.25);
  \shade[ball color=cGray!72] (2.06,3.21) ellipse (0.16 and 0.06);
  \draw[cGray!82!black, line width=0.30pt] (2.06,3.46)
    .. controls (2.44,3.49) and (2.78,3.42) .. (3.22,3.35);
  \node[draw=cGray!72!black, fill=white, line width=0.32pt, rounded corners=0.45mm,
        minimum width=0.86cm, minimum height=0.48cm,
        font=\sffamily\fontsize{4.8}{5.8}\selectfont, text=cGray!85!black] (esvm) at (3.78,3.22) {$V_s$};
  \node[small label, anchor=south] at (2.06,3.47) {ESVM probe};
  \node[small label, anchor=south] at (3.78,3.48) {$V_s$ meter};
  \draw[-{Stealth[length=2.5pt,width=1.9pt]}, cTeal!82!black, line width=0.38pt]
    (0.73,3.20)--(1.48,3.20);
  \node[small label] at (1.05,3.39) {motion stage};
  \node[font=\sffamily\fontsize{3.7}{4.4}\selectfont, text=cGray,
        anchor=east] at (4.24,2.80) {grounded substrate};
  \draw[cGray, line width=0.34pt] (2.82,2.77)--(2.82,2.69);
  \draw[cGray, line width=0.34pt] (2.69,2.69)--(2.95,2.69) (2.73,2.65)--(2.91,2.65);

  % Raw surface-potential decay.
  \draw[axis line] (0.66,1.52)--(0.66,2.49);
  \draw[axis line] (0.66,1.52)--(4.24,1.52);
  \node[small label, rotate=90] at (0.28,2.02) {$V_s(t)$};
  \node[small label, anchor=north east] at (4.21,1.44) {$t$};
  \draw[cRed!76!black, line width=0.68pt]
    (0.88,2.37) .. controls (1.36,2.04) and (1.86,1.70) .. (2.72,1.58)
                .. controls (3.20,1.52) and (3.66,1.52) .. (4.02,1.52);
  \foreach \x/\y in {1.10/2.20,1.66/1.84,2.35/1.63} {
    \shade[ball color=cRed!48] (\x,\y) circle (0.048);
    \draw[cRed!75!black, line width=0.25pt] (\x,\y) circle (0.048);
  }

  % Explicit raw-to-derived transformation.
  \draw[-{Stealth[length=2.8pt,width=2.1pt]}, cGray!72!black, line width=0.42pt]
    (3.62,1.44)--(3.62,1.28);
  \node[small label, anchor=east] at (3.50,1.35) {derive};

  % Derived trap-density distribution in energy space.
  \draw[axis line] (0.66,0.36)--(0.66,1.27);
  \draw[axis line] (0.66,0.36)--(4.24,0.36);
  \node[small label, rotate=90] at (0.28,0.82) {$g(E_t)$};
  \node[small label, anchor=north east] at (4.21,0.29) {$E_t$};
  \path[draw=cBlue!82!black, fill=cBlue!18, line width=0.58pt]
    (0.92,0.36) .. controls (1.20,0.36) and (1.27,0.88) .. (1.56,0.88)
                .. controls (1.85,0.88) and (1.92,0.36) .. (2.20,0.36) -- cycle;
  \path[draw=cRed!82!black, fill=cRed!18, line width=0.66pt]
    (2.22,0.36) .. controls (2.53,0.36) and (2.58,1.12) .. (3.02,1.12)
                .. controls (3.46,1.12) and (3.51,0.36) .. (3.84,0.36) -- cycle;
  \foreach \x/\y in {1.32/0.70,1.56/0.88,1.80/0.70} {
    \shade[ball color=cBlue!62] (\x,\y) circle (0.043);
  }
  \foreach \x/\y in {2.70/0.78,3.02/1.12,3.34/0.78} {
    \shade[ball color=cRed!62] (\x,\y) circle (0.043);
  }
  \draw[cGray!68!black, line width=0.30pt] (1.56,1.20)--(3.02,1.20);
  \draw[cGray!68!black, line width=0.30pt] (1.56,1.15)--(1.56,1.25) (3.02,1.15)--(3.02,1.25);
  \node[small label, anchor=south] at (2.29,1.24) {$\tau_d$};
  \node[small label, text=cBlue!82!black, anchor=north] at (1.56,0.31) {shallow};
  \node[small label, text=cRed!82!black, anchor=north] at (3.02,0.31) {deep};
\end{scope}

% Panel F
\begin{scope}[shift={(10.00,0)}]
  \draw[panel frame] (0,0) rectangle (4.75,4.35);
  \node[panel letter, anchor=west] at (0.28,3.97) {F};
  \node[panel title] at (0.82,3.97) {Floating Coulomb response};

  \PanelFFloatingCantilever{f}{(2.28,3.08)}
  \node[small label, anchor=east] at (1.72,3.48) {mechanical clamp};
  \draw[leader] (1.76,3.40)--(2.24,3.38);
  \node[small label, anchor=east, align=right] at (1.415,1.54) {floating polymer\\cantilever};
  \draw[leader] (1.31,1.38)--(1.58,1.62);
  \node[small label, anchor=west, align=left] at (3.83,1.55) {driven\\electrode};
  \draw[leader] (4.04,1.60)--(3.70,1.68);
  \node[small label, anchor=west, align=left] at (4.06,3.50) {grounded\\source return};
  \draw[leader] (4.02,3.45)--(3.98,3.41);
  \draw[-{Stealth[length=3.5pt,width=2.7pt]}, cRed!90!black, line width=0.82pt]
    (1.48,1.76)--(0.68,1.76);
  \node[small label, text=cRed!88!black, anchor=south] at (1.08,1.84) {Coulomb force};
  \node[small label, text=cRed!88!black, anchor=west] at (0.46,0.53) {trapped charge};
  \draw[leader, cRed!88!black] (1.42,0.63)--(1.20,1.12);
\end{scope}

\end{tikzpicture}
\end{document}
